\section{Relation to earlier theories and the retained mechanisms}\label{sec:comparison}
The projection identity \eqref{eq:fixedcp}, belief-state reduction, concavity of finite-horizon Bayes risk, conditional expectation and positive-kernel comparison are standard tools. The contribution here is the exact pooling identity \eqref{eq:poolidentity} combined with integrated curvature, actual posterior-density converses, adaptive action selection and a complete autonomous state count. Theorem~\ref{thm:pooling} is not a claim that finite-memory control, belief quantization or active sensing is new.

\subsection{Finite-memory control and quantization}
Kara and Y\"uksel~\cite{kara} approximate a partially observed control problem by a finite window of actions and reports. Under filter stability they establish policy-error bounds in terms of window length; their finite-action, finite-observation construction admits general hidden state spaces. Our reports are continuous and a retained label is not a stored observation window. Lemma~\ref{lem:aggregation} computes its true conditional hidden distribution under the finite observer's own law. The hypotheses are different: posterior smoothing and an interior carrier replace the stability hypothesis, while the conclusion matches a terminal squared-risk memory exponent and counts the internal phase. Neither theorem contains the other.

Approximate information states~\cite{ais} control planning loss through reward and transition approximation. Belief covering~\cite{zhang} makes geometric complexity relevant to POMDP planning, and filter quantization~\cite{pages} provides recursive finite approximations of conditional distributions. Here a pointwise approximation error does not suffice: it would give a first-order bound at a nonsmooth policy switch. The zero conditional first moment in \eqref{eq:barycenter}, together with integration of the curvature measure, gives a second-order bound even at such switches. Moreover, the representatives are not required to approximate the full-history posterior pathwise. Static quantization~\cite{graf} supplies the classical terminal problem; it does not itself give this causal realization.

\subsection{Observation-adaptive sensing}
Adaptive submodularity~\cite{golovin} gives approximation guarantees for adaptive stochastic utility optimization and coverage. It concerns the value or test cost of a policy, not the number of distinct states in a persistent controller. No submodularity of the present squared audit risk is asserted. For coverage comparisons one must distinguish the corrected author version's additional assumptions and squared-logarithmic guarantee from the original, subsequently corrected theorem.

Active sequential hypothesis testing~\cite{naghshvar} balances acquisition and decision costs through sequential policies and information bounds. Controlled sensing with Markovian observations~\cite{nitinawarat} permits causal sensing, nonuniform costs and asymptotic testing guarantees. Those tasks optimize stopping and identification rather than a fixed-horizon categorical forecast under a hard internal-state budget. Theorem~\ref{thm:pooling} instead fixes the physical horizon, compares against the optimal full-history sensor policy, and supplies a memory converse in the same loss. It does not replace their sequential or asymptotic optimality assertions.

Noisy optimal decision trees~\cite{jia} optimize test cost for identifying a hypothesis, including persistent noisy outcomes. Their objective and noise model differ from squared hidden-state prediction. A policy tree and a finite persistent observer are not synonymous: every internal node is part of the autonomous state count unless a separate implementation merges it. Supplement T's separated-refinement proof explicitly pays the entire tree; the present observer merges incoming edges by conditional barycenters and pays the entire phase-labelled state set. Existence of that programme is not an efficient decision-tree synthesis theorem.

Active feature acquisition also changes the next observation with the already acquired values. The acquisition-conditioned oracle~\cite{valancius} constructs nongreedy acquisition policies for predictive tasks. The time-varying evaluation theory of von Kleist et al.~\cite{kleist} addresses distribution shifts induced by changing acquisition policies, using explicit causal assumptions. We likewise distinguish the actual law under each policy, but assume known calibrated kernels and prove a finite-state geometric risk law; we do not estimate its policy value from retrospectively acquired data. In particular \eqref{eq:poolidentity} uses the law of the implemented controller, not data distributed according to another exploration policy.

\subsection{Prediction, comparison and causal information}
Predictive-state representations~\cite{psr} and causal-state constructions~\cite{shalizi} motivate the future-test equivalence relation. Appendix~\ref{sec:quotient} states the additional measurable realization conditions used here; a formal quotient is not silently assumed standard Borel. Blackwell comparison and statistical deficiency~\cite{blackwell,lecam} supply the classical decision-theoretic basis for Proposition~\ref{prop:morphism}. The additional causal assertion requires the simulator's actual internal memory, strategy lift, clock and joint audit law.

The retained continuation-cut bounds in Supplement T are finite functional coding problems with decoder side information. Their relationship with Wyner--Ziv and coding for computing~\cite{wynerziv,orlitsky} is retained there. They expose intermediate information that terminal quantization can miss. The present pooling identity is instead an upper implementation theorem for adaptive controlled hidden-state experiments, with a direct actual-mass lower. It is not a proof that arbitrary continuation widths always generate a recursion.

\subsection{What has and has not been unified}
The separated-refinement completion theorem in Supplement T requires positive branch mass, global covering and a response-support gap; it proves a horizon-uniform two-budget comparison. Theorem~\ref{thm:pooling} removes that support gap on a positive noisy sensor class and gives a full-history adaptive comparison, but retains the explicit horizon factor in \eqref{eq:autrate}. The serial continuation and sparse-filter theorems in the integral supplements remain independent sufficient mechanisms, with all their proofs. Their fixed-controller or fixed-acquisition assertions have not been promoted to unrestricted optimal control.

A unifying decision principle is now explicit: forgetting replaces a posterior by a conditional barycenter, and its exact loss is a Jensen defect of the optimal executable continuation value. Separated geometry controls a related loss by actual-mass tree frontiers; smooth acquired laws control it by integrated curvature. This principle does not yet imply a near-necessary characterization, horizon-uniform matching under overlapping noise, unknown-kernel learning, or optimal programme, workspace and planning complexity. Those conclusions require further proofs, not reinterpretation of the present inequalities.
